%============================================================================%
% Young-Hsiang TC mortality regressions with ERA5-Land wind
%
% Figures/tables are produced by
%   adhoc/wind_tc_validation/02_young_hsiang_regressions_era5.R
% Compile from this directory:  latexmk -pdf wind_tc_validation.tex
%============================================================================%
\documentclass[10pt, aspectratio=169]{beamer}

\usetheme{default}
\usecolortheme{dove}
\setbeamertemplate{navigation symbols}{}
\setbeamertemplate{caption}{\insertcaption}
\setbeamerfont{title}{series=\bfseries}
\setbeamerfont{frametitle}{series=\bfseries}
\setbeamercolor{frametitle}{fg=black}

\usepackage{booktabs}
\usepackage{graphicx}
\usepackage{amsmath}
\usepackage[T1]{fontenc}
\usepackage{lmodern}

\graphicspath{{../figures/wind_tc_validation/}}

% \yhCumRef, \corrTC, ... written by the R script
\input{../tables/wind_tc_validation/facts.tex}
\input{../tables/wind_tc_validation/diag_facts.tex}
\input{../tables/wind_tc_validation/sweep_facts.tex}

\AtBeginSection[]{%
  \begin{frame}[plain]
    \vfill
    \centering
    {\usebeamerfont{title}\insertsectionhead}
    \vfill
  \end{frame}
}

\title{Can ERA5-Land Wind Stand In for LICRICE TC Wind?}
\subtitle{Replicating Young \& Hsiang, and swapping the exposure measure}
\date{\today}

\begin{document}

\begin{frame}[plain]
  \titlepage
\end{frame}

\begin{frame}{Roadmap}
  \begin{enumerate}
    \item \textbf{Setup} --- the panel, and the specification we are matching.
    \item \textbf{Replication} --- does the Young--Hsiang headline reproduce? \emph{Yes.}
    \item \textbf{Comparing the two wind measures} --- descriptive comparison of \texttt{maxs} and ERA5-Land wind.
    \item \textbf{Results}
    \item \textbf{Why the ERA5 coefficients are larger} --- diagnostics: scaling, masking, season and state placebos, and leave-one-out.
  \end{enumerate}
  \vfill
  \footnotesize
  \textbf{Summary.} 
  The replication reproduces the published estimate. ERA5-Land wind is associated with mortality in the cyclone season, but does not measure the same events set as \texttt{maxs}.
\end{frame}

%============================================================================%
\section{Setup}
%============================================================================%

\begin{frame}{Purpose}
  \begin{enumerate}
    \item \textbf{Replicate} the Young--Hsiang (2024) main result: a 20-year
          distributed lag of tropical-cyclone wind on all-cause monthly
          mortality.
    \item \textbf{Swap} their LICRICE wind measure (\texttt{maxs}) for the
          ERA5-Land wind we merged onto the same panel, and run the two
          against each other.
  \end{enumerate}
  \vfill
  \textbf{Panel.} U.S. state $\times$ month, \yhStates{} states,
  \yhYearLo--\yhYearHi{} in estimation (lags reach back to 1930).
  Outcome is all-cause deaths per 100,000 per month.
  \vfill
  \footnotesize
  Panel built by the merge documented in
  \texttt{2\_analysis/wind\_tc\_validation/README.md}; original code at
  \texttt{Global-Policy-Lab/young\_hsiang\_tc\_mortality}.
\end{frame}

\begin{frame}{The specification}
  Young and Hsiang estimate, for state $s$ in month-of-sample $\tau$,
  \begin{align*}
    y_{s\tau} \;=\;& \sum_{j=0}^{240} \beta_j \, W_{s,\tau-j}
      \;+\; \underbrace{\mu_{sm}}_{\substack{\text{state}\,\times\\\text{cal.\ month}}}
      \;+\; \underbrace{\delta_\tau}_{\substack{\text{month-of-}\\\text{sample}}} \\[0.6em]
      &+\; \underbrace{f_s(\tau)}_{\substack{\text{state 8th-order}\\\text{polynomial in }\tau}}
      \;+\; \underbrace{g_s(T_{s\tau})}_{\substack{\text{state quadratic}\\\text{in temperature}}}
      \;+\; \varepsilon_{s\tau}
  \end{align*}
  plus a state-by-calendar-month linear trend. The object of interest is the
  \textbf{cumulative} effect $B(h) = \sum_{j=0}^{h}\beta_j$; they report it at
  $h = \yhLagRef$ months.
  \vfill
  \footnotesize
  Their \texttt{areg} uses iid standard errors, so we match that (caveats on the
  replication slide). We estimate in \texttt{fixest}, absorbing the same column space
  (241 wind lags, $\approx$2{,}000 absorbed parameters, \yhNobs{} observations).
\end{frame}

%============================================================================%
\section{Replication}
%============================================================================%

\begin{frame}{Replication: the Young--Hsiang main result}
  \centering
  \includegraphics[width=0.72\linewidth]{replication_main.png}
\end{frame}

\begin{frame}{Replication: caveats}
\begin{itemize}
  \item \texttt{tavg} is missing for 2013--2015, so the temperature control truncates estimation at \yhYearHi. This is a property of Young and Hsiang's own data: their published \texttt{.dta} has the same gap, so their regressions truncate at \yhYearHi{} too.
  \item With 36 states and 241 lag coefficients the state-clustered VCOV is rank-deficient (cluster rank $\le 36$), so clustered SEs are not reported here --- Young and Hsiang use randomization inference rather than a cluster correction. We have not reproduced that step. \textbf{We use `iid' SEs here.}
\end{itemize}
\end{frame}

%============================================================================%
\section{Comparing the two wind measures}
%============================================================================%

\begin{frame}{ERA5-Land against \texttt{maxs}, state-month by state-month}
  \centering
  \includegraphics[width=0.92\linewidth]{scatter_maxs_era5.png}
  \vfill
  \raggedright
  \footnotesize
  Black dots in (b) are binned means. The fitted slope is \olsSlope: a 1 m/s
  increase in TC wind shows up as about \olsSlope{} m/s of ERA5-Land wind. In
  (a), the spike of points on the vertical axis is the \maxsShrZero\% of
  state-months with no tropical cyclone --- ERA5-Land wind there ranges over
  essentially its whole support.
\end{frame}

\begin{frame}{The same two series through time}
  \centering
  \includegraphics[height=0.84\textheight]{timeseries.png}
\end{frame}

\begin{frame}{How the two measures differ}
  \begin{itemize}
    \item \textbf{Attribution.} \texttt{maxs} is LICRICE wind modelled from
          cyclone tracks, and is zero in \maxsShrZero\% of state-months.
          ERA5-Land is the monthly maximum of daily maximum reanalysis wind,
          whatever the cause; its minimum is \eraMin{} m/s, mean \eraMean.
    \item \textbf{Scale.} Regressing ERA5-Land on \texttt{maxs} within TC months
          gives a slope of \olsSlope.
    \item \textbf{Seasonality.} By monthly mean, raw ERA5-Land peaks in
          \peakRawMean{} and \texttt{maxs} in September. Across all state-months
          $r = \corrAll$; within TC months $r = \corrTC$.
    \item \textbf{Composition.} Of the 50 highest ERA5-Land wind months,
          \topEraTC{} are TC months. The remainder are mostly winter months in
          Nebraska and Iowa, states with almost no tropical cyclones.
  \end{itemize}
\end{frame}

\begin{frame}{Seasonality, and separation of TC months}
  \centering
  \includegraphics[height=0.52\textheight]{diagnostics.png}
  \vfill
  \raggedright
  \footnotesize
  \textbf{Top.} \texttt{maxs} and the exceedance measure peak in cyclone season on
  every statistic; raw ERA5-Land does not --- its 99th percentile peaks in
  \peakRawPtile{} and its mean in \peakRawMean, both lowest in July--August. Its
  single \emph{maximum} does fall in \peakRawMax, but that is one observation
  (Rhode Island, 1954) on a month whose 99th percentile is the second lowest of
  the year: \topHundredAug{} of the 100 windiest state-months are in August, and
  \topHundredSeason{} of 100 in June--November against 50 if timing were flat.
  \textbf{Bottom.} Within TC states the densities largely overlap (\tcMonthMean{}
  vs \nonTcMonthMean{} m/s), so the raw level does not separate TC months.
\end{frame}

%============================================================================%
\section{Swapping in ERA5-Land}
%============================================================================%

\begin{frame}{The three ERA5-Land wind measures}
  All are state-month, in m/s. $w_{s\tau}$ is the ERA5-Land monthly maximum of
  daily maximum wind, averaged over state $s$.
  \begin{enumerate}
    \item \textbf{Area-weighted} $w_{s\tau}$ --- the variable as delivered by the
          merge, grid cells averaged by area.
    \item \textbf{Population-weighted} --- same, averaged by LandScan population
          instead. Reported as a robustness; it behaves like (1) throughout.
    \item \textbf{Exceedance}, written \texttt{w\_exc}:
      \[
        w^{\text{exc}}_{s\tau} \;=\; \max\!\left(0,\;\; w_{s\tau} - q^{95}_{s,m(\tau)}\right)
      \]
      where $q^{95}_{s,m}$ is the 95th percentile of ERA5-Land wind in state $s$
      and \emph{calendar month} $m$, over \eraYearMin--\yhYearHi.
  \end{enumerate}
\end{frame}

\begin{frame}{Sample restrictions}
  ERA5-Land wind starts in \eraYearMin. A 240-lag model needs 20 years of wind
  history before the first usable outcome, so the ERA5 regression begins in
  \eraSampleLo, 20 years later than the replication.
  \vfill
  Run different samples: 
  \begin{enumerate}
    \item \textbf{\texttt{maxs}, full sample.} The replication from the previous
          section, estimated \yhYearLo--\yhYearHi.
    \item \textbf{\texttt{maxs}, ERA5-feasible sample.} The \emph{same variable
          and the same specification}, re-estimated only on
          \eraSampleLo--\yhYearHi. Nothing about the measurement changes --- only
          the sample length. This is the control.
    \item \textbf{ERA5-Land area-weighted wind}, on that same
          \eraSampleLo--\yhYearHi{} sample. This is the actual swap.
  \end{enumerate}

\end{frame}


\begin{frame}{1-year window}
  \yearSampleLo--\yhYearHi, \yearNobs{}
  obs
  \vfill
  \centering
  \includegraphics[width=\linewidth]{head_to_head_12.png}
  \vfill
  \raggedright
  \footnotesize
  All series are signed at 12 months ($t$ = \yrMaxsT, \yrEraRawT, \yrExcT;
  masked, \yrRawMaskT{} and \yrExcMaskT). Masked to cyclone months, raw ERA5
  falls from \yrRawSD{} to \yrRawMaskSD, below \texttt{maxs} (\yrMaxsSD), while
  exceedance changes little (\yrExcSD{} $\to$ \yrExcMaskSD). Masked series use
  their \emph{parent's} SD.
\end{frame}

\begin{frame}{5-year window}
  \shortSampleLo--\yhYearHi{} for every measure.
  \vfill
  \centering
  \includegraphics[width=\linewidth]{head_to_head.png}
  \vfill
  \raggedright
  \footnotesize
  The masked series remain signed at this horizon ($t = \shortRawMaskT$ raw,
  $t = \shortExcMaskT$ exceedance).
\end{frame}

\begin{frame}{Cumulative effect by lag window: first \sweepShortMax{} lags}
  L = \sweepShortStep, 2$\times$\sweepShortStep, \dots, \sweepShortMax, on the
  identical samples, in SD units on a common vertical scale. Bars are 95\%
  intervals.
  \vfill
  \centering
  \includegraphics[width=0.55\linewidth]{lag_sweep_short.png}
  \vfill
  \raggedright
  \footnotesize
  On its natural sample \texttt{maxs} is near zero to about 50 lags, then rises
  to \sweepShortAhundred{} at \sweepShortMax{} ($t = \sweepShortAhundredT$). On
  the fixed \eraSampleLo{} sample the ERA5 series reach \sweepShortBraw{} (raw)
  and \sweepShortBexc{} (exceedance), with intervals widening in $L$.
\end{frame}

\begin{frame}{Cumulative effect by lag window}
L = \sweepStep, 2$\times$\sweepStep, \dots, 240 lags. 
  \vfill
  \centering
  \includegraphics[width=0.80\linewidth]{lag_sweep.png}
\end{frame}

\begin{frame}{All lags included}
  \centering
  \includegraphics[width=0.60\linewidth]{sensitivity_240lags.png}
  \vfill
  \raggedright
  \small
  \textbf{(1) vs (2).} Same variable and specification, shorter sample: the
  cumulative effect moves from \yhCumEnd{} to \sensMaxsEnd, a factor of five with
  no change in measurement.
  \vspace{0.4em}

  \textbf{(2) vs (3).} Identical sample (\eraSampleLo--\yhYearHi, \sensNobs{}
  observations) and identical specification; only the wind variable differs.
\end{frame}

\begin{frame}{Horse race: both measures entered jointly}
  \centering
  \includegraphics[width=0.78\linewidth]{horse_race.png}
  \vfill
  \raggedright
  \footnotesize
  Entered jointly, the ERA5 exceedance measure remains signed ($t = \raceExcT$)
  while \texttt{maxs} attenuates ($t = \raceMaxsT$).
\end{frame}

%============================================================================%
\section{Why the ERA5 coefficients are larger}
%============================================================================%

\begin{frame}{Four candidate explanations}

  \begin{enumerate}
    \item \textbf{Scaling.} ``Per SD'' may not be equivalent for a sparse and
          an always-on series.
    \item \textbf{Non-cyclone weather.} ERA5-Land may capture other damaging
          weather associated with mortality.
    \item \textbf{Season or climate.} The association may be seasonal rather
          than storm-specific.
    \item \textbf{Influential units.} A few states or years may drive it.
  \end{enumerate}
  \vfill
  Everything that follows is on the \diagLags-month window, \diagNobs{}
  observations, produced by
  \texttt{03\_era5\_diagnostics.R}.
\end{frame}

\begin{frame}{(1) Scaling}
  What ``one standard deviation'' means for each measure:
  \vfill
  \centering
  \small
  \input{../tables/wind_tc_validation/diag_scaling.tex}
  \vfill
  \raggedright
  \footnotesize
  The last column is the size of a 1-SD move relative to a typical event in that variable. \\
  It is of similar magnitude for \texttt{maxs} and the exceedance measure, so the per-SD comparison is not the source of the ordering. \\
  Raw ERA5 wind is the outlier on sparsity: it is positive in every state-month, so its SD describes routine variation rather than events.
\end{frame}

\begin{frame}{(2)--(3) Where the ERA5 signal sits}
  \centering
  \includegraphics[width=0.72\linewidth]{diag_decomposition.png}
  \vfill
  \raggedright
  \footnotesize
  Each split is entered jointly, so the pair partitions the parent variable.
  Units are one SD \emph{of the parent measure}. \texttt{maxs} (green) is absent
  from the first block because it defines that split (its non-TC half is
  identically zero), and from the no-TC states, which hold 9 TC months between
  them. Its Dec--May arm rests on 76 events, hence the wide interval.
\end{frame}

\begin{frame}{Diagnostic results}
  \centering
  \tiny
  \input{../tables/wind_tc_validation/diag_results.tex}
  \vfill
  \raggedright
  \footnotesize
  Cumulative effect at \diagLags{} months, deaths per 100,000 per m/s.
  Each ``split'' block enters both components in one regression, so they
  partition the parent variable.
\end{frame}

\begin{frame}{Richer climate controls}
  The baseline already carries Young and Hsiang's temperature control --- a
  state-specific quadratic in \emph{contemporaneous} monthly mean temperature.
  \vfill
  \centering
  \includegraphics[width=0.62\linewidth]{diag_controls.png}
  \vfill
  \raggedright
\end{frame}

\begin{frame}{On controlling for rainfall}
  We cannot, with this panel. Its only climate fields are \texttt{tavg} and the
  ERA5-Land wind columns --- there is no precipitation variable in it. Could be added. 
\end{frame}

\begin{frame}{(4) Leave one state out, leave one year out}
  \centering
  \includegraphics[width=0.86\linewidth]{diag_jackknife.png}
  \vfill
  \raggedright
  \footnotesize
  \texttt{maxs} is the fragile one: dropping a single unit moves it by up to
  \jkMaxsRange\%, and its $t$ falls to \jkMaxsMinT{} (dropping
  \maxsWorstYear; \maxsWorstState{} is the most influential state). The
  exceedance measure moves at most \jkExcRange\% and never drops below
  $t = \jkExcMinT$.
\end{frame}

\begin{frame}{Summary of the diagnostics}
  \begin{enumerate}
    \item The ERA5 result survives every placebo.
    \item \textbf{But it is not measuring the same event set as \texttt{maxs}.}
          Roughly half of it sits in June--November months that Young and Hsiang
          code as having \emph{no} tropical cyclone. Candidates: storms below
          their attribution threshold, near-misses whose wind field still
          reaches the state, and other severe convective weather in the same
          season.
    \item \textbf{\texttt{maxs} is a narrower, noisier instrument.} It is zero in
          \maxsShrZero\% of months, so it is identified off relatively few
          events --- hence the jackknife fragility. ERA5-Land uses more of the panel.
  \end{enumerate}
\end{frame}


\end{document}
